% Tema 1: Deduccion natural en logica proposicional
% Diapositivas de clase. Compilar con: pdflatex tema1_deduccion_natural.tex (x2)
%
% Version para repartir (sin pausas): ver tema1_deduccion_natural_handout.tex
\documentclass[11pt,aspectratio=169]{beamer}

\usetheme{metropolis}
\metroset{sectionpage=progressbar, numbering=fraction, progressbar=frametitle,
          block=fill}

\usepackage[T1]{fontenc}
\usepackage[utf8]{inputenc}
\usepackage[spanish,es-noquoting]{babel}
\usepackage{alltt}
\usepackage{booktabs}
\usepackage{amssymb}
\usepackage{stmaryrd}

\definecolor{azul}{HTML}{1F4E79}
\definecolor{verde}{HTML}{2E7D32}
\definecolor{granate}{HTML}{8E2A2A}
\setbeamercolor{frametitle}{bg=azul}
\setbeamercolor{progress bar}{fg=verde}
\setbeamercolor{alerted text}{fg=granate}
\setbeamercolor{block title}{bg=azul!12, fg=azul}
\setbeamercolor{block body}{bg=azul!5}

% --- simbolos de Isabelle, tal como los muestra jEdit -----------------
\newcommand{\sAnd}{\(\wedge\)}
\newcommand{\sOr}{\(\vee\)}
\newcommand{\sNot}{\(\neg\)}
\newcommand{\sImp}{\(\longrightarrow\)}
\newcommand{\sIff}{\(\longleftrightarrow\)}
\newcommand{\sMeta}{\(\Longrightarrow\)}
\newcommand{\sTurn}{\(\vdash\)}
\newcommand{\sLb}{\(\llbracket\)}
\newcommand{\sRb}{\(\rrbracket\)}

% --- coloreado del codigo, siguiendo la salida LaTeX de Isabelle ------
% (mismos colores del libro de 2018: comandos violeta, cadenas rojo
%  ladrillo, comentarios verde)
\definecolor{isaCmd}{HTML}{A020F0}   % comandos de Isar
\definecolor{isaStr}{HTML}{B22222}   % cadenas "..."
\definecolor{isaCmt}{HTML}{008000}   % comentarios (* ... *)
\definecolor{resalte}{HTML}{FFE9A8}  % fondo de las reglas

\newcommand{\kw}[1]{\textcolor{isaCmd}{\textbf{#1}}}
\newcommand{\str}[1]{\textcolor{isaStr}{#1}}
\newcommand{\cmt}[1]{\textcolor{isaCmt}{\textit{#1}}}

% --- resaltado de reglas de deduccion ---------------------------------
% Fondo en vez de color de letra: asi no compite con la paleta anterior
% y la regla queda inconfundible al proyectar.
% \Rn: nombre de regla en prosa y en los listados de reglas
\newcommand{\Rn}[1]{\textcolor{granate}{\textbf{#1}}}
% \R: dentro de una demostracion, donde hay que localizar la regla entre
% los demas tokens. Fondo en vez de color de letra, para no competir con
% la paleta de Isabelle.
\newcommand{\R}[1]{{\setlength{\fboxsep}{1pt}%
  \colorbox{resalte}{\ttfamily\bfseries\textcolor{granate}{#1}}}}
% Nota: el bloque de reglas NO puede ser una macro. El contenido va en
% alltt, y si se pasa como argumento de macro TeX consume los saltos de
% linea antes de que alltt los active, y las reglas se fusionan en un
% parrafo. Se usa el entorno block con overlay: \begin{block}<2->{...}
\newcommand{\nota}[1]{{\small\itshape\color{gray} #1}}

\title{Deducción natural en lógica proposicional}
\subtitle{Tema 1 --- Razonamiento Automático}
\author{Fabián Fernando Serrano Suárez}
\institute{Ciencias de la Computación}
\date{Isabelle/HOL}

\begin{document}

\maketitle

% =====================================================================
\section{Cómo se lee una demostración}

\begin{frame}[fragile]{Anatomía de una demostración en Isar}
\begin{alltt}\footnotesize
\kw{lemma} ejemplo_1_1:
  \kw{assumes} 1: \str{"p \sAnd{} q"} \kw{and}
          2: \str{"r"}
  \kw{shows} \str{"q \sAnd{} r"}
\kw{proof} -
  \kw{have} 3: \str{"q"} \kw{using} 1 \kw{by} (rule \R{conjunct2})
  \kw{show} 4: \str{"q \sAnd{} r"} \kw{using} 3 2 \kw{by} (rule \R{conjI})
\kw{qed}
\end{alltt}
  
  \vspace{-0.4em}
  \begin{columns}[t]
    \begin{column}{0.48\textwidth}
      \small
      \texttt{{\bf{assumes}}} \quad para indicar las hipótesis\\
      \texttt{{\bf{and}}} \quad separa las hipótesis\\
      \texttt{{\bf{shows}}} \quad la conclusión\\
      \texttt{{\bf{proof -}}} \quad prueba sin método por defecto
    \end{column}
    \begin{column}{0.48\textwidth}
      \small
      \texttt{{\bf{have}}} \quad un paso intermedio\\
      \texttt{{\bf{using}}} \quad qué hechos usa el paso\\
      \texttt{\bf{by (rule R)}} \quad con qué regla\\
      \texttt{{\bf{show}}} / \texttt{{\bf{qed}}} \quad la conclusión y el cierre.
    \end{column}
  \end{columns}
  \vspace{0.6em}
  
  \nota{Cada paso dice explícitamente \alert{qué hechos se utilizan} y \alert{con qué reglas}.}
\end{frame}

\begin{frame}{Dos niveles de flecha}
  \begin{center}
  \begin{tabular}{lll}
    \toprule
    \bfseries Símbolo & \bfseries Se lee & \bfseries Nivel\\
    \midrule
    
    $\longrightarrow$ & «si \ldots{} entonces» & fórmula del \alert{objeto}\\
    
    $\Longrightarrow$ & «de \ldots{} se deduce» & regla del \alert{meta-nivel}\\
    
    $\llbracket A;\, B \rrbracket \Longrightarrow C$ & «de $A$ y $B$ se deduce $C$» & regla con dos premisas\\
    
    $\vdash$ & «\ldots{} demuestra \ldots» & el enunciado del ejercicio\\
    \bottomrule
  \end{tabular}
  \end{center}
  \vfill
  
  \begin{center}
    \alert{$p \longrightarrow q$ es una fórmula que puede ser verdadera o
    falsa. $P \Longrightarrow Q$ es una regla de inferencia.}
  \end{center}
\end{frame}

% =====================================================================
\section{Conjunción}

\begin{frame}[fragile]{Las reglas de la conjunción}
  \begin{block}{Reglas de deducción usadas}
\begin{alltt}\footnotesize
\Rn{conjI}:      \sLb{}P; Q\sRb{} \sMeta{} P \sAnd{} Q
\Rn{conjunct1}: P \sAnd{} Q \sMeta{} P
\Rn{conjunct2}: P \sAnd{} Q \sMeta{} Q
\end{alltt}
\end{block}
  
  \begin{itemize}
    \item \Rn{conjI} \emph{introduce} la conjunción: de $P$ y $Q$, se
      concluye $P \wedge Q$.
    
    \item \Rn{conjunct1} y \Rn{conjunct2} la \emph{eliminan}: de
      $P \wedge Q$ se deduce cada parte.
  \end{itemize}
  \vfill
  
  \nota{Toda regla de deducción natural viene en pareja:
  una de introducción y una de eliminación por conectivo.}
\end{frame}

\begin{frame}[fragile]{Ejemplo 1 --- \texttt{p \sAnd{} q, r \sTurn{} q \sAnd{} r}}
  \begin{columns}[T]
    \begin{column}{0.54\textwidth}
\begin{alltt}\footnotesize
\kw{lemma} ejemplo_1_1:
  \kw{assumes} 1: \str{"p \sAnd{} q"} \kw{and}
          2: \str{"r"}
  \kw{shows} \str{"q \sAnd{} r"}
\kw{proof} -
  \kw{have} 3: \str{"q"}
    \kw{using} 1 \kw{by} (rule \R{conjunct2})
  \kw{show} 4: \str{"q \sAnd{} r"}
    \kw{using} 3 2 \kw{by} (rule \R{conjI})
\kw{qed}
\end{alltt}
    \end{column}
    \begin{column}{0.44\textwidth}
      \begin{block}<2->{Reglas de deducción usadas}
\begin{alltt}\scriptsize
\Rn{conjunct2}: P \sAnd{} Q \sMeta{} Q
\Rn{conjI}: \sLb{}P; Q\sRb{} \sMeta{} P \sAnd{} Q
\end{alltt}
\end{block}
      \footnotesize
      \begin{itemize}
        \item<3-> \textbf{Paso 3:} de la hipótesis 1
          (\texttt{p \sAnd{} q}) se deduce \texttt{q}.
        \item<4-> \textbf{Paso 4:} junto \texttt{q} (paso 3) con
          \texttt{r} (hipótesis 2).
      \end{itemize}
    \end{column}
  \end{columns}
\end{frame}

\begin{frame}[fragile]{La misma prueba, más corta}
  \textbf{Dejando la regla implícita} con \texttt{..}\ --- Isabelle la
  encuentra:
\begin{alltt}\footnotesize
\kw{proof} -
  \kw{have} \str{"q"} \kw{using} assms(1) ..
  \kw{thus} \str{"q \sAnd{} r"} \kw{using} assms(2) ..
\kw{qed}
\end{alltt}
  
  \textbf{Automática:}
\begin{alltt}\footnotesize
\kw{lemma} ejemplo_1_5: \str{"\sLb{}p \sAnd{} q; r\sRb{} \sMeta{} q \sAnd{} r"}
\kw{by} auto
\end{alltt}
  \vspace{0.4em}
  
  \begin{center}
    \alert{Siete líneas $\longrightarrow$ dos.}\\
  \end{center}
\end{frame}

% =====================================================================
\section{Doble negación}

\begin{frame}[fragile]{Ejemplo 2 --- \texttt{p, \sNot\sNot(q \sAnd{} r) \sTurn{} \sNot\sNot p \sAnd{} r}}
  \begin{columns}[T]
    \begin{column}{0.54\textwidth}
\begin{alltt}\footnotesize
\kw{lemma} ejemplo_2_1:
  \kw{assumes} 1: \str{"p"} \kw{and}
          2: \str{"\sNot\sNot(q \sAnd{} r)"}
  \kw{shows} \str{"\sNot\sNot p \sAnd{} r"}
\kw{proof} -
  \kw{have} 3: \str{"\sNot\sNot p"}
    \kw{using} 1 \kw{by} (rule \R{notnotI})
  \kw{have} 4: \str{"q \sAnd{} r"}
    \kw{using} 2 \kw{by} (rule \R{notnotD})
  \kw{have} 5: \str{"r"}
    \kw{using} 4 \kw{by} (rule \R{conjunct2})
  \kw{show} 6: \str{"\sNot\sNot p \sAnd{} r"}
    \kw{using} 3 5 \kw{by} (rule \R{conjI})
\kw{qed}
\end{alltt}
    \end{column}
    \begin{column}{0.44\textwidth}
      \begin{block}<2->{Reglas de deducción usadas}
\begin{alltt}\scriptsize
\Rn{notnotI}: P \sMeta{} \sNot\sNot P
\Rn{notnotD}: \sNot\sNot P \sMeta{} P
\end{alltt}
\end{block}
      \footnotesize
      \begin{itemize}
        \item<3-> Cuatro pasos, cada uno con su regla. Se lee como el
          razonamiento que uno haría en papel.
      \end{itemize}
    \end{column}
  \end{columns}
\end{frame}

% =====================================================================
\section{Condicional}

\begin{frame}[fragile]{Ejemplo 3 --- el modus ponens}
  Demostrar \quad
  \texttt{\sNot p \sAnd{} q,\; \sNot p \sAnd{} q \sImp{} r \sOr{} \sNot p
  \;\sTurn{}\; r \sOr{} \sNot p}
  \vspace{0.5em}
  
\begin{alltt}\footnotesize
\kw{lemma} ejemplo_3_1:
  \kw{assumes} 1: \str{"\sNot p \sAnd{} q"} \kw{and}
          2: \str{"\sNot p \sAnd{} q \sImp{} r \sOr{} \sNot p"}
  \kw{shows}      \str{"r \sOr{} \sNot p"}
\kw{proof} -
  \kw{show} \str{"r \sOr{} \sNot p"} \kw{using} 2 1 \kw{by} (rule \R{mp})
\kw{qed}
\end{alltt}
  
  \begin{block}{Reglas de deducción usadas}
\begin{alltt}\footnotesize
\Rn{mp}: \sLb{}P \sImp{} Q; P\sRb{} \sMeta{} Q
\end{alltt}
\end{block}
  
  \vspace{0.2em}
  \nota{Un solo paso. El orden de \texttt{using 2 1} importa: primero el
  condicional, después el antecedente.}
\end{frame}

\begin{frame}[fragile]{Ejemplo 8 --- introducir un condicional}
  Demostrar \quad \texttt{p \sImp{} q \;\sTurn{}\; \sNot q \sImp{} \sNot p}
  \vspace{0.4em}
  
  \begin{columns}[T]
    \begin{column}{0.54\textwidth}
\begin{alltt}\footnotesize
\kw{lemma} ejemplo_8_1:
  \kw{assumes} 1: \str{"p \sImp{} q"}
  \kw{shows} \str{"\sNot q \sImp{} \sNot p"}
\kw{proof} -
  \{ \kw{assume} 2: \str{"\sNot q"}
    \kw{have} \str{"\sNot p"} \kw{using} 1 2 \kw{by} (rule \R{mt}) \}
  \kw{thus} \str{"\sNot q \sImp{} \sNot p"} \kw{by} (rule \R{impI})
\kw{qed}
\end{alltt}
    \end{column}
    \begin{column}{0.44\textwidth}
      \begin{block}<3->{Reglas de deducción usadas}
\begin{alltt}\scriptsize
\Rn{impI}: (P \sMeta{} Q) \sMeta{} P \sImp{} Q
\Rn{mt}: \sLb{}P \sImp{} Q; \sNot Q\sRb{} \sMeta{} \sNot P
\end{alltt}
\end{block}
      \footnotesize
      \begin{itemize}
        \item<4-> Las llaves \texttt{\{ \ldots{} \}} abren un
          \alert{bloque}: dentro vale la suposición \texttt{\sNot q}.
        \item<5-> \Rn{impI} convierte ese bloque ---«suponiendo
          \texttt{\sNot q} se obtiene \texttt{\sNot p}»--- en el
          condicional.
      \end{itemize}
    \end{column}
  \end{columns}
  \vspace{0.3em}
  \onslide<6->{\nota{Dentro del bloque razonamos con \(\Longrightarrow\); el
  resultado es una fórmula con \(\longrightarrow\).}}
\end{frame}

% =====================================================================
\section{Disyunción}

\begin{frame}[fragile]{Las reglas de la disyunción}
  \begin{block}{Reglas de deducción usadas}
\begin{alltt}\footnotesize
\Rn{disjI1}: P \sMeta{} P \sOr{} Q
\Rn{disjI2}: Q \sMeta{} P \sOr{} Q
\Rn{disjE}:  \sLb{}P \sOr{} Q; P \sMeta{} R; Q \sMeta{} R\sRb{} \sMeta{} R
\end{alltt}
\end{block}
  
  \begin{itemize}
    \item Introducir es fácil: basta una de las dos partes.
    
    \item \alert{Eliminar exige razonar por casos}: hay que llegar a la
      misma conclusión $R$ suponiendo $P$ y suponiendo $Q$.
  \end{itemize}
  \vfill
  
  \nota{\Rn{disjE} es la primera regla del tema con tres premisas. Se
  reconoce el patrón: «o pasa esto o pasa aquello, y en ambos casos\ldots».}
\end{frame}

\begin{frame}[fragile]{Ejemplo 12 --- \texttt{p \sOr{} q \sTurn{} q \sOr{} p}}
  \begin{columns}[T]
    \begin{column}{0.54\textwidth}
\begin{alltt}\footnotesize
\kw{lemma} ejemplo_12_1:
  \kw{assumes} \str{"p \sOr{} q"}
  \kw{shows} \str{"q \sOr{} p"}
\kw{proof} -
  \kw{have} \str{"p \sOr{} q"} \kw{using} assms \kw{by} \kw{this}
  \kw{moreover}
  \{ \kw{assume} 2: \str{"p"}
    \kw{have} \str{"q \sOr{} p"} \kw{using} 2 \kw{by} (rule \R{disjI2}) \}
  \kw{moreover}
  \{ \kw{assume} 3: \str{"q"}
    \kw{have} \str{"q \sOr{} p"} \kw{using} 3 \kw{by} (rule \R{disjI1}) \}
  \kw{ultimately} \kw{show} \str{"q \sOr{} p"} \kw{by} (rule \R{disjE})
\kw{qed}
\end{alltt}
    \end{column}
    \begin{column}{0.44\textwidth}
      \footnotesize
      \begin{itemize}
        \item<2-> \texttt{moreover} va \alert{acumulando} hechos;
          \texttt{ultimately} se los entrega todos juntos a la regla.
        \item<3-> Los tres hechos acumulados son exactamente las tres
          premisas que pide \Rn{disjE}:
          \begin{enumerate}
            \item<4-> \texttt{p \sOr{} q}
            \item<4-> de \texttt{p} sale \texttt{q \sOr{} p}
            \item<4-> de \texttt{q} sale \texttt{q \sOr{} p}
          \end{enumerate}
      \end{itemize}
    \end{column}
  \end{columns}
\end{frame}

% =====================================================================
\section{Negación}

\begin{frame}[fragile]{Las reglas de la negación}
  \begin{block}{Reglas de deducción usadas}
\begin{alltt}\footnotesize
\Rn{notI}:   (P \sMeta{} False) \sMeta{} \sNot P
\Rn{notE}:   \sLb\sNot P; P\sRb{} \sMeta{} R
\Rn{FalseE}: False \sMeta{} P
\end{alltt}
\end{block}
  
  \begin{itemize}
    \item \Rn{notI}: para probar \texttt{\sNot P}, suponga $P$ y llegue a
      contradicción.
    
    \item \Rn{notE}: de $P$ y $\neg P$ sale \alert{cualquier cosa}.
  \end{itemize}
  \vfill
  
  \nota{\Rn{FalseE} de falso se deduce cualquier fórmula.}
\end{frame}

\begin{frame}[fragile]{Ejemplo 15 --- \texttt{\sNot p \sOr{} q \sTurn{} p \sImp{} q}}
  \begin{columns}[T]
    \begin{column}{0.54\textwidth}
\begin{alltt}\footnotesize
\kw{lemma} ejemplo_15_1:
  \kw{assumes} 1: \str{"\sNot p \sOr{} q"}
  \kw{shows} \str{"p \sImp{} q"}
\kw{proof} (rule \R{impI})
  \kw{assume} 2: \str{"p"}
  \kw{note} 1
  \kw{thus} \str{"q"}
  \kw{proof} (rule \R{disjE})
    \{ \kw{assume} 3: \str{"\sNot p"}
      \kw{show} \str{"q"} \kw{using} 3 2 \kw{by} (rule \R{notE}) \}
  \kw{next}
    \{ \kw{assume} 4: \str{"q"}
      \kw{show} \str{"q"} \kw{using} 4 \kw{by} \kw{this} \}
  \kw{qed}
\kw{qed}
\end{alltt}
    \end{column}
    \begin{column}{0.44\textwidth}
      \footnotesize
      \begin{itemize}
        \item<2-> \textbf{Hacia atrás:} \texttt{proof (rule \Rn{impI})}
          dice con qué regla se cerrará \alert{antes} de probar nada.
        \item<3-> Dentro, \Rn{disjE} abre los dos casos, separados por
          \texttt{next}.
        \item<4-> \textbf{Caso \texttt{\sNot p}:} junto con \texttt{p} da
          contradicción, y \Rn{notE} produce la \texttt{q}.
        \item<5-> \textbf{Caso \texttt{q}:} ya la tenemos.
      \end{itemize}
    \end{column}
  \end{columns}
\end{frame}

% =====================================================================
\section{Bicondicional y reglas clásicas}

\begin{frame}[fragile]{Bicondicional}
  \begin{block}{Reglas de deducción usadas}
\begin{alltt}\footnotesize
\Rn{iffI}:  \sLb{}P \sMeta{} Q; Q \sMeta{} P\sRb{} \sMeta{} P \sIff{} Q
\Rn{iffD1}: \sLb{}Q \sIff{} P; Q\sRb{} \sMeta{} P
\Rn{iffD2}: \sLb{}P \sIff{} Q; Q\sRb{} \sMeta{} P
\end{alltt}
\end{block}
  
  \vfill
  \begin{center}
    Probar una equivalencia es probar \alert{las dos implicaciones}.
  \end{center}
\end{frame}

\begin{frame}[fragile]{Ejemplo 24 --- razonar por contradicción}
  Demostrar \quad \texttt{\sNot p,\; p \sOr{} q \;\sTurn{}\; q}
  \vspace{0.4em}
  
  \begin{columns}[T]
    \begin{column}{0.54\textwidth}
\begin{alltt}\footnotesize
\kw{lemma} ejemplo_24_1:
  \kw{assumes} \str{"\sNot p"}
          \str{"p \sOr{} q"}
  \kw{shows} \str{"q"}
\kw{using} `p \sOr{} q`
\kw{proof} (rule \R{disjE})
  \kw{assume} \str{"p"}
  \kw{with} assms(1) \kw{show} \str{"q"} \kw{by} \R{contradiction}
\kw{next}
  \kw{assume} \str{"q"}
  \kw{thus} \str{"q"} \kw{by} \kw{assumption}
\kw{qed}
\end{alltt}
    \end{column}
    \begin{column}{0.44\textwidth}
      \begin{block}<3->{Reglas de deducción usadas}
\begin{alltt}\scriptsize
\Rn{ccontr}:
  (\sNot P \sMeta{} False) \sMeta{} P
\end{alltt}
\end{block}
      \footnotesize
      \begin{itemize}
        \item<4-> \texttt{by contradiction} cierra el caso: tenemos
          \texttt{p} y \texttt{\sNot p} a la vez.
        \item<5-> \alert{\Rn{ccontr} es clásica}: no vale en lógica
          intuicionista. Es la que permite «supongamos que no\ldots».
      \end{itemize}
    \end{column}
  \end{columns}
\end{frame}

% =====================================================================
\section{Resumen}

\begin{frame}[fragile]{Todas las reglas del tema}
  \begin{center}
  \small
  \begin{tabular}{lll}
    \toprule
    \bfseries Conectivo & \bfseries Introducción & \bfseries Eliminación\\
    \midrule
    
    \(\wedge\) & \Rn{conjI} & \Rn{conjunct1}, \Rn{conjunct2}, \Rn{conjE}\\
    
    \(\vee\) & \Rn{disjI1}, \Rn{disjI2} & \Rn{disjE}\\
    
    \(\longrightarrow\) & \Rn{impI} & \Rn{mp}, \Rn{impE}\\
    
    \(\neg\) & \Rn{notI} & \Rn{notE}\\
    
    \(\longleftrightarrow\) & \Rn{iffI} & \Rn{iffD1}, \Rn{iffD2}\\
    
    \(\bot\) & --- & \Rn{FalseE}\\
    \bottomrule
  \end{tabular}
  \end{center}
  \vfill
  
  \small
  \textbf{Derivadas:} \Rn{mt}, \Rn{notnotI}, \Rn{notnotD}\qquad
  \textbf{Clásicas:} \Rn{ccontr}, \Rn{classical}, \Rn{excluded\_middle}
  \vfill
\end{frame}

\begin{frame}
  \begin{itemize}
    \item Cada paso de una prueba dice \alert{de qué hechos se deduce} y
      \alert{con qué reglas}.
    
    \item Cada conectivo trae su pareja de reglas: introducción y
      eliminación.
    
    \item Distinga \(\longrightarrow\) (fórmula) de \(\Longrightarrow\)
      (regla). Es el error más común del tema.
    
    \item Escriba primero la versión detallada. La automática viene
      después, y se entiende porque ya sabe qué está automatizando.
  \end{itemize}
  \vfill
\end{frame}

\end{document}
